% Integration slot for the supplied phase-aware dense-reward manuscript.
\subsubsection{Reward in the LIBERO Component Study}\label{sec:dense_method}
The LIBERO implementation represents supported pick-and-place behavior with five ordered phases: reach, grasp, lift, move, and place. Each environment maintains its own tracker. If $p_t\in[0,1]$ denotes active-phase progress and $B_t$ the best progress previously reached in that phase, the rewarded increment is $\Delta p_t^+=\max(p_t-B_t,0)$. The baseline resets on phase advancement. Revisiting an already reached state within a phase earns no additional progress credit.

Let $c_t$ count newly completed phases and $y_t$ denote terminal verifier success. The additive reward is
\begin{equation}
r_t=\lambda y_t+\operatorname{clip}(w_p\Delta p_t^++w_c c_t,-C,C).
\label{eq:dense}
\end{equation}
Only shaping is clipped; the original terminal reward is retained without a second terminal term. Setting $w_p=0$ gives phase-only shaping: intermediate credit occurs at phase boundaries. With $w_p>0$, new progress can earn credit before the next boundary. Thus phase awareness and reward density are separate choices. Both variants remember achievements, as CDPR milestones do, but neither specifies CDPR's stage-local credit or prefix protection. \pending{Insert exact phase predicates, progress functions, transition/reset logic, and chunk-to-reward aggregation from source code.}

Rollout selection is a separate mechanism. A group of $n$ trajectories with terminal outcomes $Y_i$ is retained when $L\le n^{-1}\sum_iY_i\le U$. Bounds excluding zero reject all-failure groups even if dense returns differ. Dense feedback therefore does not automatically remove this filtering bottleneck. Both the reward law and the acceptance rule must be reported when interpreting update throughput.
